\chapter{Conformation by convexity}

\eyebrow{shape and muscling, from the profile}

\vspace{4pt}
Conformation describes the shape and muscle development of the carcass. The
grading norm defines it through the \textbf{convexity of the profile}: full muscle
pushes the outline outward (convex); poor muscle, with bone apparent, sinks it
inward (concave). Apex transcribes exactly this rule into a measurement.

\section{The method}

The \textbf{integral-invariant convexity} measure (Manay et al.\ 2004; Pottmann
et al.\ 2009) computes, at each point of the carcass contour, how much of a small
disk falls inside the silhouette. That area carries the signed curvature without
differentiating, so it is robust to mask noise, and yields a value:

\begin{center}
\ttfamily $c > 0$ \rm convex (muscled) \quad\textbullet\quad
\ttfamily $c = 0$ \rm straight \quad\textbullet\quad
\ttfamily $c < 0$ \rm concave (bone apparent)
\end{center}

It is \textbf{analytical} (no training, no weights), \textbf{invariant} to
rotation, translation, scale and lighting, and uses only the segmented silhouette,
with no colour or texture. The value is averaged per anatomical region: \emph{leg}
(perna), \emph{loin} (lombo) and \emph{shoulder} (paleta), with the hung pose
(hock on top) fixing the axis.

\section{Reading the convexity map}

The \textbf{Conformation} tab (Figure~\ref{fig:analysis-conformation}) shows, per
carcass, the contour coloured by convexity over a dimmed photo: \textbf{blue where
convex} (muscled), \textbf{red where concave}. Each region's mean convexity is
annotated, and beneath the map the leg/loin/shoulder indices are listed with an
estimated grade pill.

\shot{analysis-conformation.jpeg}{Analysis $\rightarrow$ Conformation. The convexity
map reproduces the anatomy the grading norm describes: the legs read blue
(convex, muscled), the medial notch between them reads red (a recess). Per-region
indices and an estimated grade appear below each map.}{analysis-conformation}

Notice how the map is consistent with anatomy without any label: the legs are
blue, the notch between them is red, the loin is near-neutral. This
correspondence is verifiable by eye.

\begin{caution}[title=The estimated grade is not validated]
The convexity map and the per-region indices are \textbf{objective measurements}.
The \emph{estimated grade} derived from them (shown with an \tele{est.} pill in
amber) is \textbf{not validated}. The project's oracle experiment showed that
conformation is not reliably recoverable from a single 2D image at this sample
size (Station~13). Use the grade as an indicative reference, to be validated as
the paired dataset grows.
\end{caution}

\section{Where the numbers appear}

Conformation runs automatically during Analysis (when the background is removed).
For each analysed carcass, the Inspector shows the leg / loin / shoulder
convexities, the consolidated index, and the estimated grade
(Figure~\ref{fig:carcass-images}, right). The CSV export includes all of them.
